% 第 34 章骨架：导引框已定，正文待写。写作计划见 docs/PLAN.md“第 34 章”，写法见 docs/STYLE.md。
\chapter{战争财政、功能财政与布雷顿森林的设计}\label{ch:warfinance}

\begin{chapterguide}
\guidequestion{战争证明政府“付得起”，和平时期也一样吗？战后的国际货币体系应该怎样设计？}
\begin{guidegoals}
\item 用凯恩斯《如何筹措战费》（1940）分析战时的通胀缺口与强制储蓄；
\item 概括勒纳（1943）功能财政的规则与卡莱茨基（1943）的政治经济学批评；
\item 比较美英、苏联与纳粹德国的战时经济组织；
\item 比较凯恩斯方案与怀特方案，说明 1944 年体系的最终设计。
\end{guidegoals}
\guideprereq{\chref{ch:islm}\enspace\nameref{ch:islm}；\chref{ch:alternatives}\enspace\nameref{ch:alternatives}}
\end{chapterguide}

\chapterstub
